\chapter{Grafos}


\section{Representaciones comunes}

\subsection*{Lista de adyacencia}

\begin{lstlisting}[caption=Lista de adyacencia]
	vector<int> G[MAXN]; // grafo sin peso
	
	void agregar_arista(int u, int v) {
		G[u].push_back(v);
		G[v].push_back(u); // si es no dirigido
	}
\end{lstlisting}

\subsection*{Matriz de adyacencia}

\begin{lstlisting}[caption=Matriz de adyacencia]
	bool mat[MAXN][MAXN];
	
	void agregar_arista(int u, int v) {
		mat[u][v] = true;
		mat[v][u] = true; // si es no dirigido
	}
\end{lstlisting}

\section{Recorridos clásicos}

\subsection{DFS (Depth-First Search)}

\begin{lstlisting}[caption=DFS no dirigido]
	bool vis[MAXN];
	vector<int> G[MAXN];
	
	void dfs(int u) {
		vis[u] = true;
		for (int v : G[u]) {
			if (!vis[v]) dfs(v);
		}
	}
\end{lstlisting}

\subsection{BFS (Breadth-First Search)}

\begin{lstlisting}[caption=BFS con cálculo de distancias]
	int dist[MAXN];
	bool vis[MAXN];
	vector<int> G[MAXN];
	
	void bfs(int s) {
		queue<int> q;
		q.push(s);
		vis[s] = true;
		dist[s] = 0;
		
		while (!q.empty()) {
			int u = q.front(); q.pop();
			for (int v : G[u]) {
				if (!vis[v]) {
					vis[v] = true;
					dist[v] = dist[u] + 1;
					q.push(v);
				}
			}
		}
	}
\end{lstlisting}

\section{Componentes conexas}

\begin{lstlisting}[caption=Contar componentes conexas]
	int comp = 0;
	
	for (int i = 0; i < n; ++i) {
		if (!vis[i]) {
			dfs(i);
			comp++;
		}
	}
\end{lstlisting}

\section{Dijkstra (caminos mínimos)}

\begin{lstlisting}[caption=Dijkstra con priority queue]
	typedef pair<int, int> pii;
	vector<pii> G[MAXN]; // {peso, destino}
	int dist[MAXN];
	
	void dijkstra(int s) {
		priority_queue<pii, vector<pii>, greater<pii>> pq;
		fill(dist, dist + MAXN, INF);
		dist[s] = 0;
		pq.push({0, s});
		
		while (!pq.empty()) {
			auto [d, u] = pq.top(); pq.pop();
			if (d > dist[u]) continue;
			for (auto [w, v] : G[u]) {
				if (dist[u] + w < dist[v]) {
					dist[v] = dist[u] + w;
					pq.push({dist[v], v});
				}
			}
		}
	}
\end{lstlisting}

\section{Bellman-Ford}

\begin{lstlisting}[caption=Bellman-Ford para grafos con pesos negativos]
	vector<tuple<int, int, int>> edges; // {u, v, peso}
	int dist[MAXN];
	
	bool bellman(int s, int n) {
		fill(dist, dist + MAXN, INF);
		dist[s] = 0;
		
		for (int i = 0; i < n - 1; ++i) {
			for (auto [u, v, w] : edges) {
				if (dist[u] + w < dist[v]) {
					dist[v] = dist[u] + w;
				}
			}
		}
		
		// Detectar ciclos negativos
		for (auto [u, v, w] : edges) {
			if (dist[u] + w < dist[v]) return true;
		}
		
		return false;
	}
\end{lstlisting}

\section{Problemas clásicos}

\subsection*{A - Número de componentes conexas}

\textbf{Estado:} nodos visitados.  
\textbf{Recorrido:} DFS o BFS.  
\textbf{Complejidad:} \( O(n + m) \)

\subsection*{B - Camino más corto (no ponderado)}

\textbf{Idea:} usar BFS desde el nodo origen y guardar las distancias.

\subsection*{C - Dijkstra / Bellman-Ford}

\textbf{Cuando usar cada uno:}

- Dijkstra → aristas no negativas.
- Bellman-Ford → permite pesos negativos.

\section{¿Cómo verificar si un grafo es una línea (path)?}



\begin{enumerate}
	\item Verificar que \( M = N - 1 \) (condición necesaria para ser un árbol).
	\item Verificar que el grafo es conexo (DFS o BFS desde cualquier nodo).
	\item Contar los grados de los nodos:
	\begin{itemize}
		\item Debe haber exactamente 2 nodos con grado 1.
		\item Todos los demás deben tener grado 2.
	\end{itemize}
\end{enumerate}

\subsection*{Código en C++}

\begin{lstlisting}[caption=Verificar si un grafo es una línea]
	vector<int> G[MAXN];
	bool vis[MAXN];
	int deg[MAXN];
	
	void dfs(int u) {
		vis[u] = true;
		for (int v : G[u]) {
			if (!vis[v]) dfs(v);
		}
	}
	
	bool es_linea(int n, int m) {
		if (m != n - 1) return false; // debe tener N-1 aristas
		
		// verificar conectividad
		fill(vis, vis + n, false);
		dfs(0);
		for (int i = 0; i < n; ++i)
		if (!vis[i]) return false;
		
		// contar grados
		int cnt_deg1 = 0;
		for (int i = 0; i < n; ++i) {
			if (deg[i] == 1) cnt_deg1++;
			else if (deg[i] != 2) return false;
		}
		
		return cnt_deg1 == 2;
	}
\end{lstlisting}

\textbf{Entrada:}
\begin{itemize}
	\item Llenar \texttt{G[i]} con las aristas del grafo.
	\item Aumentar \texttt{deg[u]} y \texttt{deg[v]} cada vez que se agrega una arista entre \( u \) y \( v \).
\end{itemize}
		
\section{Detectar ciclos en grafo no dirigido}

\begin{enumerate}
	\item Usar DFS manteniendo el nodo padre
	\item Si encuentras un nodo ya visitado que no es el padre, hay ciclo
\end{enumerate}

\subsection*{Código en C++}

\begin{lstlisting}[caption=Detectar ciclos en grafo no dirigido]
	vector<int> G[MAXN];
	bool vis[MAXN];
	
	bool hasCycleUndirected(int u, int parent) {
		vis[u] = true;
		for (int v : G[u]) {
			if (!vis[v]) {
				if (hasCycleUndirected(v, u)) 
				return true;
			}
			else if (v != parent) 
			return true;
		}
		return false;
	}
	
	bool tieneCicloNoDirigido(int n) {
		fill(vis, vis + n, false);
		for (int i = 0; i < n; ++i) {
			if (!vis[i] && hasCycleUndirected(i, -1))
			return true;
		}
		return false;
	}
\end{lstlisting}

\section{Detectar ciclos en grafo dirigido}

\begin{enumerate}
	\item Usar DFS con colores (0: no visitado, 1: en proceso, 2: visitado)
	\item Si encuentras un nodo "en proceso" durante el DFS, hay ciclo
\end{enumerate}

\subsection*{Código en C++}

\begin{lstlisting}[caption=Detectar ciclos en grafo dirigido]
	vector<int> G[MAXN];
	int color[MAXN]; // 0: blanco (no visitado), 1: gris (en proceso), 2: negro (visitado)
	
	bool hasCycleDirected(int u) {
		color[u] = 1; // gris
		for (int v : G[u]) {
			if (color[v] == 0) {
				if (hasCycleDirected(v)) 
				return true;
			} 
			else if (color[v] == 1) 
			return true;
		}
		color[u] = 2; // negro
		return false;
	}
	
	bool tieneCicloDirigido(int n) {
		fill(color, color + n, 0);
		for (int i = 0; i < n; ++i) {
			if (color[i] == 0 && hasCycleDirected(i))
			return true;
		}
		return false;
	}
\end{lstlisting}

\section{Verificar si hay ciclo impar}

\begin{enumerate}
	\item Intentar colorear el grafo con 2 colores (bipartición)
	\item Si no es posible, contiene ciclo impar
\end{enumerate}

\subsection*{Código en C++}

\begin{lstlisting}[caption=Detectar ciclo impar]
	vector<int> G[MAXN];
	int col[MAXN]; // 0 o 1
	
	bool tieneCicloImpar(int n) {
		queue<int> q;
		fill(col, col + n, -1);
		
		col[0] = 0;
		q.push(0);
		
		while (!q.empty()) {
			int u = q.front(); q.pop();
			for (int v : G[u]) {
				if (col[v] == -1) {
					col[v] = 1 - col[u];
					q.push(v);
				}
				else if (col[v] == col[u]) {
					return true; // Ciclo impar encontrado
				}
			}
		}
		return false;
	}
\end{lstlisting}

\section{Verificar si el grafo es bipartito}

\begin{enumerate}
	\item Igual que detectar ciclo impar, pero devolviendo el resultado inverso
	\item Si se puede colorear con 2 colores, es bipartito
\end{enumerate}

\subsection*{Código en C++}

\begin{lstlisting}[caption=Verificar si grafo es bipartito]
	vector<int> G[MAXN];
	int col[MAXN]; // 0 o 1
	
	bool esBipartito(int n) {
		queue<int> q;
		fill(col, col + n, -1);
		
		col[0] = 0;
		q.push(0);
		
		while (!q.empty()) {
			int u = q.front(); q.pop();
			for (int v : G[u]) {
				if (col[v] == -1) {
					col[v] = 1 - col[u];
					q.push(v);
				}
				else if (col[v] == col[u]) {
					return false;
				}
			}
		}
		return true;
	}
\end{lstlisting}

\section{Distancia mínima desde un nodo}

\begin{enumerate}
	\item Usar BFS para calcular distancias desde el nodo fuente
	\item Mantener arreglo de distancias inicializado en -1
\end{enumerate}

\subsection*{Código en C++}

\begin{lstlisting}[caption=Distancia mínima con BFS]
	vector<int> G[MAXN];
	int dist[MAXN];
	
	void bfs(int s, int n) {
		queue<int> q;
		fill(dist, dist + n, -1);
		
		dist[s] = 0;
		q.push(s);
		
		while (!q.empty()) {
			int u = q.front(); q.pop();
			for (int v : G[u]) {
				if (dist[v] == -1) {
					dist[v] = dist[u] + 1;
					q.push(v);
				}
			}
		}
	}
\end{lstlisting}

\section{Verificar si el grafo es un árbol}

\begin{enumerate}
	\item Verificar que tiene exactamente N-1 aristas
	\item Verificar que es conexo
	\item Verificar que no tiene ciclos
\end{enumerate}

\subsection*{Código en C++}

\begin{lstlisting}[caption=Verificar si grafo es árbol]
	vector<int> G[MAXN];
	bool vis[MAXN];
	
	bool hasCycle(int u, int parent) {
		vis[u] = true;
		for (int v : G[u]) {
			if (!vis[v]) {
				if (hasCycle(v, u)) return true;
			}
			else if (v != parent) return true;
		}
		return false;
	}
	
	bool esArbol(int n, int m) {
		if (m != n - 1) return false;
		
		// Verificar conexidad
		fill(vis, vis + n, false);
		hasCycle(0, -1);
		for (int i = 0; i < n; ++i)
		if (!vis[i]) return false;
		
		// Verificar que no tiene ciclos
		fill(vis, vis + n, false);
		if (hasCycle(0, -1)) return false;
		
		return true;
	}
\end{lstlisting}

\section{Número de hojas en un árbol}

\begin{enumerate}
	\item Contar nodos con grado 1 (excepto posiblemente la raíz en árboles enraizados)
\end{enumerate}

\subsection*{Código en C++}

\begin{lstlisting}[caption=Contar hojas en un árbol]
	vector<int> G[MAXN];
	int deg[MAXN];
	
	int contarHojas(int n) {
		int hojas = 0;
		for (int i = 0; i < n; ++i) {
			if (G[i].size() == 1) // grado 1
			hojas++;
		}
		return hojas;
	}
\end{lstlisting}

\section{Consejos}

\begin{itemize}
	\item Usá listas de adyacencia para eficiencia en espacio.
	\item Siempre inicializá `vis[]` y `dist[]` antes de usarlos.
	\item Para grafos dirigidos, no agregues `v → u` si no se especifica.
	\item En competencias, muchos problemas de grafos se reducen a BFS/DFS/Dijkstra.
	\item Si el grafo tiene pesos negativos, ¡no uses Dijkstra!
\end{itemize}
